\chapter{资料来源与分析框架}

\section{资料来源}

本示例完全基于公开资料，不包含真实病历、内部访谈、保密论文或未公开附件。资料来源主要包括媒体报道、节目时间线、公开社交平台记录以及可检索的影像和杂志信息。纳入标准强调事件时间明确、来源可追溯，并能够对应到职业阶段变化。

\section{分析框架}

研究采用“阶段划分-事件编码-韧性观察”三层结构。阶段划分用于回答职业转型在何时发生明显转折；事件编码用于概括平台迁移、角色扩展与公众反馈；韧性观察则用于讨论危机发生后，公开人物如何重新建立稳定叙事。

\section{结构说明}

在模板层面，本章同时承担“方法说明”和“版式展示”的功能。通过分章节拆分正文、保持统一段落格式和关键词表达，可以让后续用户更容易理解如何在不改变样式的前提下，替换为自己的真实学术内容。
